\documentclass{standalone}
\usepackage{pgfplots}
\pgfplotsset{compat=1.18}

\begin{filecontents*}{\jobname.bib}
@misc{dummy,
    author = {Placeholder},
    title  = {Placeholder Reference},
    year   = {1900}
}
\end{filecontents*}

\makeatletter
\AtEndDocument{%
    \immediate\write\@auxout{\string\citation{dummy}}%
    \immediate\write\@auxout{\string\bibstyle{plain}}%
    \immediate\write\@auxout{\string\bibdata{\jobname}}%
}
\makeatother

\definecolor{tabblue}{rgb}{0.0000, 0.4470, 0.7410}
\definecolor{taborange}{rgb}{0.8500, 0.3250, 0.0980}
\definecolor{tabgreen}{rgb}{0.4660, 0.6740, 0.1880}

\begin{document}
    \begin{tikzpicture}
        % --- Unit Step Respone ---
        \begin{axis}[
                    title = {\large Unit Step },
                    % title style={font=\Large},
                    name=plot1,
                    width=8cm,
                    height=5cm,
                    xmin=0,
                    xmax=1.5,
                    xtick={0,0.5,1.0,1.5},
                    ymin=0,
                    grid=both,
                    minor grid style={dashed, gray!30},              
                    major grid style={dashed, gray!30},              
                    minor tick num=1,       
                    legend pos=south east,
                    xlabel = \(t\),
                    ylabel = {\(\omega_{b_3}\)},
                ]
            \draw[dashed, black] (axis cs:\pgfkeysvalueof{/pgfplots/xmin},1) -- (axis cs:\pgfkeysvalueof{/pgfplots/xmax},1);
            \addlegendimage{dashed, black}
            \addplot[tabblue, thick] table[x=time, y=omega_z, col sep=comma] {data/wz_step_ours.csv};
            \addplot[tabgreen, thick] table[x=time, y=omega_z, col sep=comma] {data/wz_step_gazebo.csv};
            \legend{{Reference}, {Ours}, {Gazebo}}
        \end{axis}

        % --- Variable Amplitude Response ---
        \begin{axis}[
                    title = {\large Variable Amplitude},
                    name=plot2,
                    at={(plot1.right of south east)},
                    anchor=left of south west,
                    xshift=0.5cm,
                    width=8cm,
                    height=5cm,
                    xmin=0,
                    xmax=8,
                    xtick={0,4,8},
                    ymin=-6,
                    grid=both,
                    minor grid style={dashed, gray!30},              
                    major grid style={dashed, gray!30},              
                    minor tick num=1,       
                    xlabel = \(t\),
                ]
            \addplot[dashed, black, thick] table[x=time, y=omega_z, col sep=comma] {data/ref_amp.csv};
            \addplot[tabblue, thick] table[x=time, y=omega_z, col sep=comma] {data/wz_amp_ours.csv};
            \addplot[tabgreen, thick] table[x=time, y=omega_z, col sep=comma] {data/wz_amp_gazebo.csv};
        \end{axis}

        % % --- Variable Frequency Response ---
        \begin{axis}[
                    title = {\large Variable Frequency},
                    name=plot3,
                    at={(plot2.right of south east)},
                    anchor=left of south west,
                    xshift=0.5cm,
                    width=8cm,
                    height=5cm,
                    xmin=0,
                    xmax=8,
                    xtick={0,4,8},
                    ymin=-1.2,
                    grid=both,
                    minor grid style={dashed, gray!30},              
                    major grid style={dashed, gray!30},              
                    minor tick num=1,       
                    xlabel = \(t\),
                ]
            \addplot[dashed, black, thick] table[x=time, y=omega_z, col sep=comma] {data/ref_freq.csv};
            \addplot[tabblue, thick] table[x=time, y=omega_z, col sep=comma] {data/wz_freq_ours.csv};
            \addplot[tabgreen, thick] table[x=time, y=omega_z, col sep=comma] {data/wz_freq_gazebo.csv};
        \end{axis}
        \node[font=\Large\bfseries, above, yshift=10pt, xshift=12pt] at (current bounding box.north) {Yaw Response};
    \end{tikzpicture}
\end{document}